%
% $Id: slides.tex 4228 2006-06-21 21:55:12Z jjamor $
%
%
% Compilar a .pdf con LaTeX (pdflatex)
% Es necesario instalar Beamer (paquete latex-beamer en Debian)
%

%
% Gr�ficos:
% Los gr�ficos pueden suministrarse en PNG, JPG, TIF, PDF, MPS
% Los EPS deben convertirse a PDF (usar epstopdf)
%

\documentclass{beamer}
\usetheme{Warsaw}
\usebackgroundtemplate{\includegraphics[width=\paperwidth]{format/blank-bg.png}}
\usepackage[spanish]{babel}
\usepackage[latin1]{inputenc}
\usepackage{graphics}
\usepackage{amssymb} % Simbolos matematicos

%\definecolor{libresoftgreen}{RGB}{162,190,43}
%\definecolor{libresoftblue}{RGB}{0,98,143}

%\setbeamercolor{titlelike}{bg=libresoftgreen}

%% Metadatos del PDF.
\hypersetup{
  pdftitle={Ejercicio},
  pdfauthor={Gregorio Robles y Agust�n Santos},
  pdfcreator={Master EOI},
  pdfproducer=PDFLaTeX,
  pdfsubject={Repaso.}
}
%%

\begin{document}

\title{Problema}
\subtitle{Master EOI}
\institute{ }
\author{Gregorio Robles y Agust�n Santos}
\date{\today}

\frame{
\maketitle
\begin{center}
%\includegraphics[width=6cm]{format/gsyc-urjc}
\end{center}
}


% Si el titulo o el autor se quieren acortar para los pies de p�gina
% se pueden redefinir aqu�:
%\title{Titulo corto}
%\author{Autores abreviado}


%% LICENCIA DE REDISTRIBUCION DE LAS TRANSPAS
\frame{
~
\vspace{4cm}

\begin{flushright}
{\tiny
(cc) 2007-11 Grupo de Sistemas y Comunicaciones. \\
  Algunos derechos reservados. Este trabajo se entrega bajo la
  licencia Creative Commons Attribution-ShareAlike. Para obtener la
  licencia completa, v�ase
  http://creativecommons.org/licenses/by-sa/2.1/es. Tambi\'{e}n puede
  solicitarse a Creative Commons, 559 Nathan Abbott Way, Stanford,
  California 94305, USA.
}
\end{flushright}
}
%%

\section{Tipos}

\begin{frame}
	\frametitle{PROBLEMA}
	Vamos a encontrar una palabra de \alert{cinco} letras trabajando en equipo, pero..
	\begin{itemize}
  		\item No podemos hablar entre nosotros, ni usar chat.
%    	\item �nicamente podemos utilizar VisibleTweets.com.
		\item Os vamos a dejar pistas. 
      \item Os pod�is repartir el trabajo de encontrar la respuesta a cada pista.
	\end{itemize}
\end{frame}

%% ENIAC

\begin{frame}
	\frametitle{Primera Letra}
%	Utiliza la arquitectura G que hemos visto en clase. Calcula el resultado, la representaci�n en ASCII es la respuesta.
%	\begin{itemize}
%		\item Registro vale cero y recordad que tiene 8 bits.
%		\item 00000001 00000100
%		\item 00000100 00000100
%		\item 00000100 00001000
%		\item 00000010 00000101
%    	\end{itemize}

   \begin{itemize}
     \item Tenemos el siguiente valor:
     \item 00000110 11000010
     \item Codificado con 16 bits
     \item Rango de frecuencias: entre 0 y 20 KHz
     \item Obtener el nombre de la nota musical de la frecuencia obtenida
     \item Transformar a nomenclatura de nota musical anglosajona
     \item Tomar la letra dos posisiciones m�s adelante
   \end{itemize}

\end{frame}
	
	
%% N


\begin{frame}
	\frametitle{Segunda Letra}
	Para que no sea muy dif�cil, en �sta os damos pista: Est� en la misma posici�n que en el nombre del est�ndar.
\end{frame}

%% I

\begin{frame}
	\frametitle{Tercera Pista}
	\begin{enumerate}
		\item La traducci�n latina de un libro escrito por Al-Khwarizmi dio nombre a un concepto explicado hoy. 
		\item Busca en la Wikipedia (en ingl�s) ese concepto.
		\item Descarga a tu ordenador la im�gen que hay a la derecha de la p�gina (su nombre empieza por 220px).
	\item Encuentra el n�mero de veces que puedes comprimir ese fichero sin incrementar el tama�o.
	\item Escribe ese n�mero en n�meros romanos.
    \end{enumerate}
\end{frame}

%% A

\begin{frame}
	\frametitle{Cuarta Letra}
La cuarta letra es la inicial de un programador famoso.
	\begin{itemize}
		\item Encuentra el nombre de un ordenador muy conocido de una pel�cula de ciencia ficci�n.
		\item Encuentra el nombre de una compa��a que estaba una posici�n m�s retrasada.
		\item Encuentra el a�o en el que se fund� esa compa�ia y resta 51.
		\item En ese a�o se cre� el ordenador en el que nuestro personaje desconocido hizo su programa.
    \end{itemize}
\end{frame}

%% C

\begin{frame}
	\frametitle{Quinta Letra}
N�mero romano de lo que ocupa en memoria en Kb la composici�n de un texto plano en ASCII de 64.000 caracteres (de los cuales 16.250 son con tilde), a la que se le unen 24.000 instrucciones de un procesador que puede computar 64 instrucciones diferentes, m�s 3.200 estados binarios y 20.000 n�meros en ASCII.
	\begin{itemize}
		\item (No se agrupan los bits por palabras!)
\item Debe salir una cantidad exacta.
    \end{itemize}
\end{frame}

%%

\end{document}
